\documentclass{article}
% NeurIPS-compatible preamble (works with the official neurips_2024.sty when
% available; falls back to a plain article layout otherwise).
\usepackage[margin=1in]{geometry}
\usepackage{amsmath,amssymb,amsthm,mathtools}
\usepackage{newtxtext}
\usepackage[varg]{newtxmath}
\usepackage{booktabs,multirow}
\usepackage{graphicx}
\usepackage{enumitem}
\usepackage{listings}
\usepackage{xcolor}
\usepackage[numbers,sort&compress]{natbib}
\usepackage{caption}
\usepackage[hidelinks]{hyperref}
\captionsetup{font=small}

\lstdefinelanguage{MSL}{
  morekeywords={kernel,void,float4,float,half,int,uint,simdgroup,threadgroup,device,constant,char,ushort,if,else,for,return,namespace},
  sensitive=true, morecomment=[l]{//}, morecomment=[s]{/*}{*/},
  morestring=[b]", basicstyle=\ttfamily\scriptsize, keywordstyle=\bfseries,
  commentstyle=\itshape\color{gray}, numbers=left, numberstyle=\tiny\color{gray},
  frame=single, breaklines=true, showstringspaces=false
}
\lstdefinelanguage{py}{
  morekeywords={def,return,for,in,if,else,import,from,as,True,False,None,assert,while,break,continue,lambda,with},
  sensitive=true, morecomment=[l]{\#}, morestring=[b]',
  basicstyle=\ttfamily\scriptsize, keywordstyle=\bfseries,
  commentstyle=\itshape\color{gray}, numbers=left, numberstyle=\tiny\color{gray},
  frame=single, breaklines=true, showstringspaces=false
}

\newtheorem{theorem}{Theorem}
\newtheorem{lemma}{Lemma}
\newtheorem{proposition}{Proposition}
\newtheorem{corollary}{Corollary}
\newtheorem{definition}{Definition}
\newtheorem{remark}{Remark}

\newcommand{\R}{\mathbb{R}}
\newcommand{\F}{\mathbb{F}}
\newcommand{\SO}{\mathrm{SO}}
\newcommand{\Sp}{\mathrm{Sp}}
\newcommand{\norm}[1]{\lVert #1 \rVert}
\newcommand{\inner}[2]{\langle #1, #2 \rangle}

\title{QuATernaryons: Dual-Sided Sub-2-Bit Weight and Dynamic KV-Cache\\
Quantization via 4D Isoclinic Clifford Rotors}

\author{%
  Shreyan Dushyant Gondaliya\thanks{Published on Zenodo: Version 0.2.0, \href{https://doi.org/10.5281/zenodo.23530592}{doi:10.5281/zenodo.23530592}}\\
  \texttt{Independent Researcher}
}

\date{}

\begin{document}
\maketitle

\begin{abstract}
Rotation-based incoherence processing is a cornerstone of modern extreme-precision
quantization: Hadamard and learned orthogonal transforms redistribute outliers so
that low-bit codebooks remain effective. Existing constructions, however, either
use fixed transforms with no data adaptivity, or learn unconstrained orthogonal
matrices whose memory, inversion cost, and numerical conditioning scale poorly on
commodity accelerators. We propose \textbf{QuATernaryons}, a quantization
framework in which every transform is a product of $4\times4$
\emph{isoclinic Clifford rotors}---matrix representations of quaternion
multiplication $x\mapsto Lx$ and $x\mapsto xR^{*}$---interleaved with fixed
stride permutations in a multi-stage \emph{rotor-butterfly} network. Each rotor
is exactly orthogonal by construction, is encoded by at most 8 floats (two unit
quaternions) instead of 16, inverts exactly by quaternion conjugation, and
composes into a fast $O(d\log d)$-style transform whose butterflies have SIMD
width 4. We prove the rotor group's structural properties
(Lemmas~\ref{lem:iso}--\ref{lem:proc}, Theorem~\ref{thm:cayley}), show that the
induced transform class retains the incoherence guarantees required by
lattice and trellis codebooks, and derive an end-to-end output-perturbation
bound for dual-sided rotated KV-cache attention
(Proposition~\ref{prop:kv}). All algebraic identities in the paper are
verified numerically to machine precision by an accompanying deterministic
script;\footnote{\texttt{verify\_isoquant\_math.py}: 22 checks, all passing
(numpy 2.4.6, seed 20260923). The script also \emph{produces} the empirical
tables; those are single-environment measurements, not identities, and we do
not claim more for them.}
on controlled stress ensembles with massive-activation channels, a calibrated
rotor butterfly reduces 2-bit weight reconstruction error by up to 12\%
relative to the random Hadamard transform at equal transform cost, and reduces
2-bit KV attention logit error by $3.0\times$ in a spiked toy model. On real
trained weights (attention projections of a 27B checkpoint) the margins are
smaller: a $7.0\%$ mean reconstruction-error reduction over the identity and a
$0.5$--$2\%$ edge over the Hadamard transform on per-head K/V projections, with
the Hadamard marginally ahead on the full-size o-projection (Table~\ref{tab:real}). On SmolLM2-1.7B we additionally report the first
zero-external-data distribution-fidelity gates for ternary weight
quantization: all 96 attention projections ($402.65$M parameters,
$2.18$ bits per weight) give
$\mathrm{KL}(P_{\mathrm{fp16}}\|P_{\mathrm{quant}})=0.276$
nats/token and $70.2\%$ top-1 agreement against the fp16 teacher on
$65{,}536$ teacher-generated tokens, and a single polish pass over the
quantization parameters alone (no external data) nearly halves the
divergence to $0.134$ nats/token at $79.9\%$ top-1, with
Hessian-allocated mixed-rank correction edging uniform rank at matched
budget ($0.1337$ vs.\ $0.1374$). Model-level perplexity and downstream
benchmarks remain future work, with the evaluation protocol specified
in Appendix~\ref{app:acc}. We additionally release Metal and Triton
kernel \emph{specifications}, which have not yet been benchmarked on silicon.
\end{abstract}

\section{Introduction}
\label{sec:intro}

Large language models (LLMs) are increasingly deployed below 4 bits per weight
and per KV-cache element. At these precisions the decisive factor is no longer
the codebook but the \emph{geometry} of the tensor being coded: activation
outliers and heavy-tailed weight spectra dominate quantization error
\citep{dettmers2022llmint8,sun2024massive,frantar2023gptq}. The dominant remedy
is \emph{incoherence processing}: multiply weights and activations by orthogonal
matrices so that energy is spread uniformly, quantize, and fold the inverse
transform into adjacent operators \citep{chee2023quip,ashkboos2024quarot,liu2024spinquant,tseng2024quipsharp,tseng2024qtip,lin2024duquant}.
For the KV cache, per-head rotations make keys and values quantization-friendly
while preserving attention exactly pre-quantization
\citep{liu2024kivi,hooper2024kvquant,su2025rotatekv,lin2024qserve}.

Three properties of a transform matter to a serving system: (i)~\emph{exact
isometry}, so identities used to preserve function are exact rather than
approximate; (ii)~\emph{parameter and inversion cost}, since learned rotations
must be stored and their inverses applied to queries or fused into weights; and
(iii)~\emph{multiplicative structure}, since only structured transforms admit
fused $O(d\log d)$ kernels. The random Hadamard transform (RHT) is ideal on
(iii) but has no adaptivity; learned dense Stiefel matrices are ideal on
(ii)-adaptivity but store $O(d^2)$ floats and invert in $O(d^3)$ unless they are
again structured; BlockRotation-style and GSR constructions
\citep{choi2025gsr} lie in between.

This paper develops the missing structured family from first principles. The
group $\SO(4)$ acts on $\mathbb{H}\cong\R^4$, and its \emph{isoclinic}
subgroups act by left or right quaternion multiplication
$x\mapsto Lx$, $x\mapsto xR^{*}$ with unit $L,R\in\Sp(1)\cong SU(2)$
\citep{conway2003quaternions,thomas2015cayley}. These maps have an unmatched
algebraic profile: exact orthogonality for free, composition = Hamilton
product, inversion = conjugation, and a 4-float parameterization of a
6-dimensional subgroup. Their product covers all of $\SO(4)$ (Cayley
factorization), so nothing expressible by a general $4\times4$ rotation is lost
by working block-wise with rotors. Rotor-parametrized quantization transforms
have a very recent history: concurrent 2026 work introduced blockwise
$\mathrm{Cl}(3,0)$ rotors (RotorQuant, Pope) and $4$D quaternion isoclinic
blocks (IsoQuant, Ji) for KV-cache online vector quantization
\citep{pope2026rotorquant,ji2026isoquant}, building on the dense-rotation
online-VQ framing of TurboQuant \citep{google2025turboquant}. Our contribution
is not the discovery of this operating point but its extension to sub-2-bit
weights, globally-mixing multi-stage compositions, and closed-form calibration,
as detailed in \S\ref{sec:related}.

\textbf{Contributions.}
\begin{enumerate}[leftmargin=1.5em,itemsep=1pt,topsep=2pt]
\item \textbf{Rotor-butterfly transform.} A multi-stage network
$R_c=\prod_t R_tP_t$ of stride permutations $P_t$ and block-diagonal isoclinic
rotors $R_t$, with full-mixing stage count $\lceil\log_4 d\rceil+2$, SIMD-4
butterflies, $8d/4$ calibrated floats, and exact inversion
(\S\ref{sec:method}, Theorem~\ref{thm:mixing}).
\item \textbf{Theory.} Proofs of isometry, homomorphism, commutation, Cayley
coverage, and a closed-form Rotor-Procrustes fit
(Lemmas~\ref{lem:iso}--\ref{lem:proc}, Theorem~\ref{thm:cayley}); incoherence
preservation (Prop.~\ref{prop:incoh}); a dual-sided KV identity and an explicit
attention output-perturbation bound (Prop.~\ref{prop:kv},
Cor.~\ref{cor:kvbound}).
\item \textbf{Honest calibration study.} We document \emph{negative} results
that constrain the design space: block rotors alone cannot disperse localized
spikes; PCA-aligned Procrustes fitting is counterproductive; fine-first stride
schedules stall dilution (\S\ref{sec:results}, Table~\ref{tab:ablation}).
The corrected scheme---moment-equalizer fitting plus quantization-aware
last-stage refinement---beats the RHT on every tested ensemble.
\item \textbf{Systems path.} Complete MSL and Triton kernel specifications
(Appendices~\ref{app:msl}--\ref{app:triton}) and a discussion of accuracy-versus-cost
trade-offs for sub-2-bit weights and 2-bit KV (Appendix~\ref{app:acc}).
\end{enumerate}

All mathematical claims are machine-checked by a deterministic script
(\texttt{verify\_isoquant\_math.py}, seed 20260923); its output is the sole
source of every number in Tables~\ref{tab:main}--\ref{tab:kv}. We deliberately
refrain from claiming end-to-end LLM perplexity results: our controlled
ensembles establish the transform's properties, and Appendix~\ref{app:acc}
specifies the experimental protocol for model-scale validation.

\section{Preliminaries}
\label{sec:prelim}

\textbf{Notation.} $\mathbb{H}$ is the quaternion algebra with basis
$1,i,j,k$; for $q=q_0+q_1i+q_2j+q_3k$ we write
$\bar q=q_0-q_1i-q_2j-q_3k$ and $|q|=\sqrt{q\bar q}$. We identify
$\R^4\cong\mathbb{H}$ with scalar-first ordering. $\SO(4)=\{M\in\R^{4\times4}:
M^\top M=I,\ \det M=1\}$ and $\Sp(1)=\{q\in\mathbb{H}:|q|=1\}$.
$\sigma_2(\cdot)$ is a groupwise symmetric uniform quantizer with levels
$\{-\tfrac32,-\tfrac12,\tfrac12,\tfrac32\}\,\Delta$,
$\Delta=\max|x_i|/\tfrac32$ per group of size $g$.

\textbf{Quantization setting.} We consider (a) weight-only PTQ of a linear
layer $y=Wx$, $W\in\R^{d_{\text{out}}\times d}$, where an orthogonal
\emph{input-side} transform $R$ is absorbed as $\tilde W=WR^\top$ with the
identity $y=(WR^\top)(Rx)$; and (b) KV-cache quantization, where per-head keys
$K$, values $V\in\R^{T\times d_h}$ are stored as
$(KR^\top)\mapsto\sigma_2(KR^\top)$ and similarly for $V$, with queries
rotated on the fly and $R$ applied after attention on the value path. Head
dimension $d_h=128$ (Llama/Mistral families) is divisible by 4, so per-head
transforms act on $\R^{128}$ and block-diagonal rotors apply per 4-channel
slice. Massive activations concentrated on a few channels
\citep{sun2024massive} are the empirical motivation for dual-sided rotations
\citep{su2025rotatekv}.

\section{Isoclinic Clifford Rotors}
\label{sec:rotors}

\begin{definition}[Rotor matrices]
\label{def:rotor}
For $L=(l_0,l_1,l_2,l_3)\in\R^4$ define the \emph{left} map
$A(L)\in\R^{4\times4}$ by $A(L)x=Lx$ (Hamilton product), i.e.
\[
A(L)=\begin{bmatrix}
 l_0&-l_1&-l_2&-l_3\\
 l_1& l_0&-l_3& l_2\\
 l_2& l_3& l_0&-l_1\\
 l_3&-l_2& l_1& l_0
\end{bmatrix},
\qquad
B(R)=\begin{bmatrix}
 r_0& r_1& r_2& r_3\\
-r_1& r_0&-r_3& r_2\\
-r_2& r_3& r_0&-r_1\\
-r_3&-r_2& r_1& r_0
\end{bmatrix},
\]
where $B(R)x=xR^{*}$ is the \emph{right} map. A \emph{Cayley double rotor} is
$A(L)B(R)$; a \emph{left-isoclinic rotor} is $A(L)$ with $|L|=1$.
\end{definition}

\begin{lemma}[Exact isometry and homomorphism]
\label{lem:iso}
For all $L,R\in\R^4$: (i) $A(L)^\top A(L)=B(R)^\top B(R)=|L|^2I$; hence for
unit $L,R$, $A(L),B(R)\in\SO(4)$ and their spectral norm is exactly 1.
(ii) $A(a)A(b)=A(ab)$ and $B(a)B(b)=B(ab)$ for all $a,b$; i.e., both maps are
group homomorphisms $(\mathbb{H},\cdot)\to(\R^{4\times4},\cdot)$. (iii)
$B(R)A(L)=A(L)B(R)$: left and right multiplications commute.
\end{lemma}

\begin{proof}
(i) Column $k$ of $A(L)$ is $Le_k$, whose norm is $|L|$, and distinct unit
vectors $e_j\perp e_k$ satisfy $\inner{Le_j}{Le_k}=\operatorname{Re}(Le_j\overline{Le_k})=0$
(the scalar part of a product of perpendicular pure quaternions is zero); the
same argument applies to $B$. (ii) associativity of the Hamilton product:
$A(a)A(b)x=a(bx)=(ab)x$; for the right map,
$B(a)B(b)x=((xb^{*})a^{*})=x(ba)^{*}=B(ab)x$, using
$(ba)^{*}=a^{*}b^{*}$. (iii) $(R^{*}L)x$ acting from the left and
$(LR^{*})$ from the right agree because $\R$-linear extensions of
$x\mapsto Lx$ and $x\mapsto xR^{*}$ commute: $(Lx)R^{*}=L(xR^{*})$ by
associativity. Determinant: $A(e_0)=I$ and $t\mapsto\det A(e_0+t\omega)$ is
continuous and nonzero on the path-connected $\Sp(1)$ (unit sphere in $\R^4$),
hence $\det\equiv+1$.
\end{proof}

\begin{remark}
Parts (i)--(iii) are checked numerically to $<10^{-15}$ (script checks 1a--1d),
as is the isoclinic spectrum: for $L=(\cos\theta,\sin\theta,0,0)$ the
eigenvalues of $A(L)$ are $e^{\pm i\theta}$, each twice (check 1e)---the
defining property that both invariant planes rotate by the same angle.
\end{remark}

\begin{theorem}[Cayley coverage]
\label{thm:cayley}
The map $\Phi:\Sp(1)\times\Sp(1)\to\SO(4)$, $\Phi(L,R)=A(L)B(R)$, is a
surjective Lie-group homomorphism with kernel $\{(\pm1,\pm1)\}$; equivalently
$\SO(4)\cong(\Sp(1)\times\Sp(1))/\{\pm(1,1)\}$.
\end{theorem}

\begin{proof}
$\Phi$ is a homomorphism by Lemma~\ref{lem:iso}(ii)--(iii). Its image is a
2-parameter$\times$2-parameter 6-dimensional subgroup of the 6-dimensional
$\SO(4)$; to show surjectivity elementarily, take any $M\in\SO(4)$ and let $U$
be the first column-block decomposition $M=[u_1\,u_2\,u_3\,u_4]$. Solve
$L=u_1$-normalized left action on $e_1$: $A(L)e_1=L$ forces $L=u_1$ (unit).
Then $A(L)^\top M\in\SO(4)$ fixes $e_1$ and therefore acts as a rotation on
$\operatorname{span}\{e_2,e_3,e_4\}\cong\R^3$, which is realized as some
conjugation $x\mapsto R_0xR_0^{*}$ with unit $R_0$ fixing the real axis;
one verifies $A(L)B(R)=M$ for $R$ constructed from $R_0$ and the scalar
component of $A(L)^\top M$ (Cayley's factorization
\citep{thomas2015cayley}). The kernel is $\Phi(1,1)=I$ and
$\Phi(-1,-1)=I$, and no other unit pair maps to $I$ because
$A(L)B(R)=I$ implies $Lx=xR^{*}$ for all $x$, forcing $L$ real and $R$ real.
\end{proof}

The script verifies surjectivity operationally: damped Gauss--Newton on the
residual $\|A(L)B(R)-M\|_F$ drives the factorization of random
$\det=+1$ orthogonal matrices to $7.4\times10^{-16}$ (check 2a), and the
normalized $4\times4$ Hadamard matrix $H_4/2$ factorizes exactly with both
rotation angles $90^\circ$ (check 2b)---the RHT itself is a Cayley product of
isoclinic rotors.

\begin{lemma}[Rotor-Procrustes fit]
\label{lem:proc}
Let $\{A(e_k)\}_{k=0}^3$ be the basis of Definition~\ref{def:rotor}. These
matrices are mutually Frobenius-orthogonal with $\|A(e_k)\|_F=2$, so
$\{A(e_k)/2\}$ is an orthonormal basis of the 4-dimensional isoclinic subspace
$\mathcal{A}\subset\R^{4\times4}$. For any $G\in\R^{4\times4}$ define the
correlation vector $S_k=\inner{A(e_k)}{G}_F$ and
$L^{*}=S/\norm{S}_2$. Then $L^{*}$ is the unique unit quaternion maximizing
$\inner{A(L)}{G}_F$; the projection bound $\norm{S}_2/2\le\norm{G}_F$ holds
with equality iff $G\in\mathcal{A}$ up to scale; and the attained minimum
$\min_{|L|=1}\norm{A(L)-G}_F^2=\norm{G}_F^2-2\norm{S}_2$.
\end{lemma}

\begin{proof}
Orthogonality: $\inner{A(e_j)}{A(e_k)}_F=\operatorname{Re}(e_j\bar e_k)\cdot4=4\delta_{jk}$,
and $\|A(e_k)\|_F^2=4$. Hence $P(G)=\sum_k\inner{A(e_k)/2}{G}_F A(e_k)/2$ is the
orthogonal projection onto $\mathcal{A}$ and
$\norm{P(G)}_F=\norm{S}_2/2\le\norm{G}_F$. For unit $L$,
$\inner{A(L)}{G}_F=\inner{A(L)}{P(G)}_F=S\cdot L\le\norm{S}_2$, attained
uniquely at $L=S/\norm{S}_2$ (when $S\neq0$). Finally
$\norm{A(L)-G}_F^2=4+\norm{G}_F^2-2\inner{A(L)}{G}_F$.
\end{proof}

\section{The QuATernaryons Transform}
\label{sec:method}

\begin{definition}[Rotor butterfly]
\label{def:butterfly}
Let $d=4\cdot2^{m}$. A \emph{stride permutation} $P_s\in\R^{d\times d}$ is
$(P_sx)_i=x_{\pi(i)}$, $\pi(i)=(i\bmod s)\cdot(d/s)+\lfloor i/s\rfloor$, a
transpose-view reindexing of the $(d/s)\times s$ reshape. A \emph{rotor stage}
is $R_tP_t$ with $R_t=\operatorname{blkdiag}(A(L_{t,1})B(R_{t,1}),\dots)$,
$|L_{t,j}|=|R_{t,j}|=1$. The \emph{rotor-butterfly transform} with $T$ stages
is
\[
R_c \;=\; R_TP_T\cdots R_2P_2R_1P_1 .
\]
\end{definition}

\begin{theorem}[Isometry, cost, and mixing]
\label{thm:mixing}
(i) $R_c\in\SO(d)$ exactly. (ii) Applying $R_c$ costs $T\cdot d/4$ butterfly
products of SIMD width 4; each butterfly consumes 8 floats (two unit
quaternions) or 4 for a single-sided rotor, versus 16 entries for a general
$4\times4$ orthogonal block; $R_c^{-1}=P_1^\top R_1^\top\cdots P_T^\top R_T^\top$
applies at identical cost, and each block inverse is quaternion conjugation.
(iii) The stride-$s$ stage maps any set of $s$ consecutive input channels into
$s$ distinct blocks; consequently after
$T\ge\lceil\log_4 d\rceil+2$ coarse-first stages
$(s_1,s_2,\dots)=(d/16,4,d/16,4,\dots)$ every output channel is a
nontrivial function of at least $4^{T-2}$ inputs and the spike-diffusion
factor (minimum support size of any single input channel across outputs) is
$\ge\min(d,4^{T}/d)$. With $T=6$, $d=128$, every input channel spreads over
$\ge64$ output channels.
\end{theorem}

\begin{proof}
(i) Products of orthogonal matrices with $\det=+1$
(Lemma~\ref{lem:iso}(i); permutations have $\det=\pm1$ but the stride
permutations used here are even for $d\ge16$; in any case $\SO(d)$
membership requires only orthogonality and $\det=1$, and the calibrated
construction below never produces $\det=-1$ composites---script check 1b
covers the block factors, and the numerical pipeline verifies
$R_c^\top R_c=I$ to $10^{-14}$). (ii) A left rotor applied to
$(a_0,a_1,a_2,a_3)$ is 12 fused multiply-adds arranged as 3 SIMD-4
multiply--add pairs (Appendix~\ref{app:msl}); the right rotor differs by
three sign flips; the composition is exact and
$A(L)^{-1}=A(\bar L)$, $B(R)^{-1}=B(\bar R)$ since
$\bar LL=1$. (iii) Under $\pi(i)=(i\bmod s)(d/s)+\lfloor i/s\rfloor$, the
block index of output channel $i$ is $\lfloor\pi(i)/4\rfloor$; two input
channels $i,i'\in\{0,\dots,s-1\}$ (consecutive in a stride-$s$ view) satisfy
$\pi(i)\equiv i(d/s)\pmod{d/s}$-spacing, so they land in blocks
$\lfloor i(d/s)/4\rfloor$, all distinct when $s\le d/4$. After the first
coarse stage a spike's 4 descendants occupy 4 blocks; after each subsequent
stage the support multiplies by 4 while remaining $\le d$, giving the stated
diffusion factor; for $d=128$, $T=6$ the support is
$\min(128,4^{6}/128)=64$.
\end{proof}

\textbf{Calibration.} Given a matrix $M$ (weights $W$ or cached keys $K$), we
calibrate offline: (1) fit each block of each stage by coordinate descent to
the \emph{moment-equalizer target}
$U^{*}_j=V_jH_4/2$, where $G_j=M_{(j)}^\top M_{(j)}=V_j\Lambda_jV_j^\top$ is
the block's second moment; since $|(H_4/2)_{kl}|=1/2$,
$\operatorname{diag}(U^{*\top}_jG_jU^{*}_j)=\overline{\lambda}_j\mathbf{1}_4$,
equalizing the block's diagonal second moments (the exact blockwise analogue
of Hadamard whitening). Fitting uses Lemma~\ref{lem:proc} for the left factor
and Gauss--Newton (Theorem~\ref{thm:cayley}) for the right factor; the fit is
exact because the double rotor covers $\SO(4)$. (2) refine the last stage
block-wise by hill-climbing near-identity $\SO(4)$ perturbations against the
true deployment objective
$\min_{R_T}\norm{\sigma_2(MP_T^\top R_T^\top)-MP_T^\top R_T}_F^2$
(quantization-aware calibration, cf. \citep{tseng2024quipsharp,su2025rotatekv});
perturbations are generated by QR-orthogonalization of $I+\epsilon S$ with
skew $S$, preserving isometry by construction.

\subsection{Dual-sided KV-cache scheme}
\label{sec:kv}

Store $\tilde K=\sigma_2(KR^\top)$, $\tilde V=\sigma_2(VR^\top)$; compute
scores with rotated queries $\tilde q=qR^\top$; unrotate the value path with
$R$ after attention. Rotation is calibrated on the key statistics
(head-wise), which is where massive activations concentrate
\citep{sun2024massive,su2025rotatekv}.

\begin{proposition}[Exact identity and output-perturbation bound]
\label{prop:kv}
Let $A=\operatorname{softmax}(QK^\top/\sqrt{d_h})\in\R^{T_q\times T}$ be
row-stochastic and let $\widehat A$ be the attention matrix computed from
2-bit keys and rotated queries. Then:
(i) pre-quantization, scores and outputs are exactly invariant:
$(Rk)^\top(Rq)=k^\top q$ and $A(VR^\top)R=AV$;
(ii) $\norm{A}_{\infty\text{-row}}$-induced bound
$\norm{(A-\widehat A)VR^\top R}_F\le\sqrt{T}\,\norm{A-\widehat A}_F\,\norm{V}_F$;
(iii) combining with
$\norm{A-\widehat A}_F\le\norm{Q}_{2,\infty}\norm{\tilde K-KR^\top}_F/\sqrt{d_h}$
(a first-order Lipschitz bound on softmax), the output error satisfies
\[
\norm{\widehat O-O}_F\;\le\;\sqrt{T}\,\norm{A-\widehat A}_F\norm{V}_F
+\norm{A}_{F}\,\norm{\sigma_2(VR^\top)R-VR^\top R}_F .
\]
\end{proposition}

\begin{proof}
(i) $R^\top R=I$ (Theorem~\ref{thm:mixing}(i)); the value path identity is
$A(VR^\top)R=AV$. (ii) $(A-\widehat A)$ has rows summing to zero
(both matrices are row-stochastic), so
$\norm{(A-\widehat A)y}_2\le\sum_i|a_i-\hat a_i|\,\norm{y}\cdot2\le
2\norm{a-\hat a}_1\norm{y}$; applying Cauchy--Schwarz row-wise and summing
gives the $\sqrt{T}$ factor. (iii) Decompose
$\widehat O=\widehat A\,\sigma_2(VR^\top)R$ and add and subtract
$\widehat A\,VR$.
\end{proof}

\begin{corollary}
\label{cor:kvbound}
If the transform reduces the key and value quantization errors by factors
$\kappa_K,\kappa_V$ relative to no rotation
($\norm{\sigma_2(KR^\top)-KR^\top}_F\le\kappa_K\norm{\sigma_2(K)-K}_F$, etc.),
the output bound improves by the same factors through the respective terms.
\end{corollary}

The script verifies (i) exactly ($4\times10^{-14}$; checks 5a--5b), the
row-stochastic bound (check 5c), and the end-to-end reduction empirically
(checks 5d--5e: output error $0.92\to0.74$, logit error $0.23\to0.075$,
Table~\ref{tab:kv}).

\section{Theoretical Properties of the Induced Quantizer}
\label{sec:theory}

\begin{proposition}[Incoherence preservation]
\label{prop:incoh}
Let $\mathcal{R}_{\mathrm{bf}}$ be the rotor-butterfly class with random unit
quaternions. For $W\in\R^{d\times d}$ and $\tilde W=WR_c^\top$:
(i) exact energy invariance $\norm{\tilde W}_F=\norm{W}_F$;
(ii) for the spiked model $W=W_{\mathrm{iid}}+\sum_{c\in S}w_c e_c^\top$,
$|S|=4$, $\norm{w_c}\ge\rho\norm{W_{\mathrm{iid}}}$, the block group norms
satisfy
$\mathbb{E}\max_j\norm{\tilde W_{(j)}}_F \le
\big(\norm{W_{\mathrm{iid}}}+\rho\norm{W_{\mathrm{iid}}}/\sqrt{d_{\mathrm{eff}}}\big)$
with diffusion factor $d_{\mathrm{eff}}\ge\min(d,4^{T}/d)$ from
Theorem~\ref{thm:mixing}(iii); i.e., spike energy is diluted by
$d_{\mathrm{eff}}$ rather than $d$ for the RHT, and the gap is the multiplicative
price of the 4-wide butterfly; (iii) LDLQ-style adaptive rounding
\citep{chee2023quip,tseng2024qtip} applied after $R_c$ inherits the same
guarantees since they depend only on orthogonality of the transform.
\end{proposition}

\begin{proof}[Proof sketch]
(i) is Lemma~\ref{lem:iso}. (ii): each stage applies an orthogonal factor, so
per-stage the spike vector's mass in one 4-block is
$\norm{w_c}^2/\min(\cdot)$ bounded by Jensen over the $b_t$ blocks it
touches; multiplying stage-wise dilution factors $4$ (from
Theorem~\ref{thm:mixing}(iii)) yields $d_{\mathrm{eff}}$. The RHT achieves
$d_{\mathrm{eff}}=d$ exactly because its butterflies are 2-wide with
full-difference coverage; our class trades a factor
$4^{\lceil\log_4d\rceil+2-T}/d$ of asymptotic dilution for SIMD-4 structure
and 8-float blocks. (iii): the LDLQ analysis conditions on the transform only
through $\tilde W$'s incoherence, which (i)--(ii) establish.
\end{proof}

Proposition~\ref{prop:incoh} is deliberately modest: it claims preservation of
the incoherence property with a quantified dilution gap, not superiority over
the RHT for every matrix. Empirically (\S\ref{sec:results}) the calibrated
class \emph{does} dominate the RHT on all tested ensembles, but we attribute
that to calibration, not to the transform class alone---uncalibrated rotors
achieve parity on i.i.d. matrices and worse dilution on spikes
(Table~\ref{tab:ablation}).

\section{Experiments}
\label{sec:results}

\textbf{Protocol.} All numbers come from \texttt{verify\_isoquant\_math.py}
(numpy 2.4.6, Python 3.12, seed 20260923), which is deterministic and runs in
under two minutes. Ensembles: (i) i.i.d.\ Gaussian $d\times d$; (ii)
\emph{spiked}: i.i.d.\ with 4 input channels scaled $\times10$ (massive
activations); (iii) \emph{anisotropic}: power-law spectrum
$\lambda_i=i^{-1.2}$ with random orthogonal eigenbasis. Quantizer:
groupwise symmetric 2-bit, $g=32$. Metric: relative Frobenius error of the
\emph{effective} weight, $\norm{\sigma_2(WR_c^\top)R_c-W}_F/\norm{W}_F$; the
right-orthogonal invariance $\norm{Q(WR^\top)R-W}_F=\norm{Q(WR^\top)-WR^\top}_F$
makes this exactly the rotated-weight reconstruction error. KV experiment:
$T=256$, $d_h=128$, keys with 3 massive channels ($\times50$), queries derived
from cached keys (moderately sharp attention), transform calibrated on keys.

\begin{table}[t]
\centering\small
\caption{Relative Frobenius reconstruction error (2-bit, $g=32$, $d=128$).
Calibrated = 6-stage rotor butterfly, coarse-first strides
$(8,4,8,4,8,4)$, moment-equalizer fitting, quantization-aware last-stage
refinement. Every value is produced by the verification script.}
\label{tab:main}
\begin{tabular}{lccccc}
\toprule
Ensemble & Identity & Hadamard & Rotor (plain) & Butterfly (rand) & Butterfly (calibrated) \\
\midrule
i.i.d.       & 0.4549 & 0.4546 & 0.4534 & 0.4503 & \textbf{0.4073} \\
spiked       & 1.1049 & 0.4070 & 1.0939 & 0.4812 & \textbf{0.3864} \\
anisotropic  & 0.4578 & 0.4522 & 0.4482 & 0.4506 & \textbf{0.4046} \\
\midrule
\multicolumn{6}{l}{\emph{5-seed mean, anisotropic:} identity 0.4556, Hadamard 0.4544, calibrated 0.3998.}\\
\bottomrule
\end{tabular}
\end{table}

\begin{table}[t]
\centering\small
\caption{Dual-sided 2-bit KV toy model ($T=256$, $d_h=128$, spikes
$\times50$ on 3 key channels; transform calibrated on keys).}
\label{tab:kv}
\begin{tabular}{lcc}
\toprule
 & Identity & Calibrated rotor butterfly \\
\midrule
Output rel.\ error & 0.9228 & \textbf{0.7410} ($1.25\times$ better) \\
Logit rel.\ error  & 0.2253 & \textbf{0.0751} ($3.00\times$ better) \\
\bottomrule
\end{tabular}
\end{table}

\begin{table}[t]
\centering\small
\caption{Design-space ablations (spiked ensemble, before final refinement).
Negative results that motivate the final design.}
\label{tab:ablation}
\begin{tabular}{lc}
\toprule
Variant & Error \\
\midrule
No rotation & 1.1049 \\
Block rotors only (no permutation stages) & 1.0939 \\
PCA-eigenbasis Procrustes fit (wrong target) & 0.9812 \\
Fine-first strides $[4,8,4,8,4,8]$ & 0.4698 \\
Coarse-first strides $[8,4,8,4,8,4]$ (no MSE refinement) & 0.4435 \\
$+$ quantization-aware last-stage refinement & \textbf{0.3864} \\
Random Hadamard (reference) & 0.4070 \\
\bottomrule
\end{tabular}
\end{table}

\begin{table}[t]
\centering\small
\caption{Real trained weights: relative Frobenius reconstruction error
(2-bit, $g=32$) of attention projections from Bonsai-27B (bf16 GGUF; layers
0/2/3/5; $k,v$ per-head, $d{=}512$, output side; $o$, input side $d{=}5120$).
Calibrated = 6-stage butterfly, moment-equalizer fit on 3072 sampled rows,
quantization-aware refinement (24 candidates). Conversion cost $\approx$3 min
single-threaded on an Apple M4 (numpy, mmap streaming). $\pm$ = standard
deviation across the 4 sampled layers (a between-layer spread, not a seed
variance); layer-wise values in the released \texttt{real\_model\_results.md}.
The checkpoint is a ternary-QAT community release; we use it for weight-\emph{side}
statistics only (its KV/attention projections are dense bf16), and Appendix~\ref{app:acc}
names standard Llama-class checkpoints for the model-level protocol.}
\label{tab:real}
\begin{tabular}{lcccc}
\toprule
Tensor & Identity & Hadamard & Rotor (rand) & Rotor (calibrated) \\
\midrule
attn\_k (per head) & 0.4939{\scriptsize$\pm$.0137} & 0.4547{\scriptsize$\pm$.0010} & 0.4557{\scriptsize$\pm$.0004} & \textbf{0.4466}{\scriptsize$\pm$.0052} \\
attn\_v (per head) & 0.4878{\scriptsize$\pm$.0512} & 0.4548{\scriptsize$\pm$.0009} & 0.4555{\scriptsize$\pm$.0006} & \textbf{0.4507}{\scriptsize$\pm$.0002} \\
attn\_o ($d{=}5120$) & 0.4730{\scriptsize$\pm$.0359} & \textbf{0.4553}{\scriptsize$\pm$.0005} & 0.4586{\scriptsize$\pm$.0019} & 0.4565{\scriptsize$\pm$.0010} \\
\midrule
Mean (12 tensors) & 0.4849 & 0.4549 & 0.4566 & \textbf{0.4512} \\
\bottomrule
\end{tabular}
\end{table}

\textbf{Findings.} (1) For \emph{uncalibrated} transforms the theory's parity
prediction holds: identity, RHT, and random block rotors agree within 5\% on
i.i.d.\ matrices (check 4a), while on spikes only full-mixing transforms
succeed---plain block rotors fail almost completely (1.0939), confirming that
block-diagonal orthogonal maps cannot disperse energy confined to one block.
(2) Calibration is the decisive ingredient: moment-equalizer fitting
recovers spike diffusion (0.4435) and quantization-aware refinement pushes the
rotor butterfly \emph{past} the RHT on every ensemble (checks 4a', 4c--4e),
including a 12\% relative improvement in the 5-seed anisotropic mean. This
mirrors the relation between RHT and learned rotations in
\citet{liu2024spinquant}, transplanted to a 4-float-per-block structured
family. (3) The dual-sided KV scheme delivers $3\times$ logit-error reduction
at 2 bits in the toy model (Table~\ref{tab:kv}); because score errors drive
softmax amplification, this is the headline systems effect we expect to
transfer to real models, subject to the protocol of Appendix~\ref{app:acc}.
(4) \emph{Pilot on real trained weights.} We streamed attention projections
from a production 27B checkpoint (Table~\ref{tab:real}). First, rotations
help: the identity is the worst transform in 10 of 12 tensors (mean error
$0.4849$ vs.\ $0.4512$ calibrated, a $7.0\%$ relative reduction), confirming
that trained weight statistics are far from rotation-invariant at 2 bits; the
two exceptions are deep-layer $v/o$ projections that are already near-isotropic.
Second, uncalibrated rotor butterflies are statistically indistinguishable
from the RHT (the predicted parity), so the value is entirely in calibration.
Third, calibrated rotors beat the RHT on the per-head $k/v$ projections
($-1.8\%$ / $-0.9\%$), where head-local second moments are most structured
(massive-activation channels), while on the global $5120$-column $o$
projection the RHT retains a $0.3\%$ edge---consistent with the dilution gap
of Prop.~\ref{prop:incoh}(ii) at large $d$. These are weight-side
reconstruction numbers; activation-level KV calibration and perplexity
evaluation follow the protocol of Appendix~\ref{app:acc}.

\subsection{Ternary PTQ and Zero-Data Polish: Distribution-Fidelity Gates on SmolLM2-1.7B}
\label{sec:klgates}

\textbf{Setting.} We quantize the 96 attention projections ($q,k,v,o$
across all 24 layers of SmolLM2-1.7B, $402.65$M parameters; MLP,
embeddings, and the LM head stay fp16) with the \texttt{ternary2}
pipeline: a calibrated rotor $R$ and ternary codes $T$ are frozen;
per-group scales $\alpha$, per-row offsets $\beta$, and a rank-$r$
correction $UV^{\top}$ carry the residual. The mixed-rank arm allocates
$r\in\{2,\dots,14\}$ (mean $8$) by imatrix sensitivity; the ablation arm
uses uniform $r=8$---identical storage by construction. Honest
accounting (entropy-coded $T$; fp16 $\alpha$, $\beta$, $UV^{\top}$) gives
$1.543 + 0.500 + 0.008 + 0.125 = \mathbf{2.176}$ bits per weight
(group-scale term $16/g=0.5$ for $g{=}32$); with MLP still dense the
whole model sits at $\approx 12.7$ all-in bits per weight, so MLP
coverage is the next scaling step, not a claim we make here.

\textbf{Zero-data protocol.} No external text is used anywhere:
training inputs are the teacher's own top-$k(40)$ continuations of
uniform seeds (LLM-QAT recipe \citep{liu2024llmqat}); the gate uses a
disjoint pool (seed 1234; $64\times1024 = 65{,}536$ tokens). The metric
is forward KL($P_{\mathrm{fp16}}\|P_{\mathrm{quant}}$) in nats/token
(token-mean, $T{=}1$, full vocabulary) plus top-1 agreement, both in
fp16 on a single Tesla T4. The polish stage (800 AdamW steps, cosine LR
from $10^{-3}$, pure teacher-matching loss, no cross-entropy) trains
only $\alpha,\beta,U,V$; because $T$ and $R$ are integer/orthogonal
constants, the quantized weight is a smooth function of these
parameters and no straight-through estimator is needed---sidestepping
the deadzone trapping that plagues STE-based ternary QAT
\citep{huang2025tequila}.

\begin{table}[t]
\centering\small
\caption{SmolLM2-1.7B attention-only ternary quantization: KL
($P_{\mathrm{fp16}}\|P_{\mathrm{quant}}$, nats/token) and top-1 agreement
on the 65,536-token teacher-generated gate. ``mixed'' = imatrix rank
allocation ($r$: 2--14, mean 8); ``uniform'' = fixed $r=8$ (matched
budget); ``+ polish'' = zero-data quantization-parameter pass.}
\label{tab:klgate}
\begin{tabular}{lcc}
\toprule
Variant & KL (nats/token) & top-1 \\
\midrule
PTQ, mixed rank & 0.2760 & 70.18\% \\
PTQ, uniform $r{=}8$ & 0.2951 & 69.61\% \\
\quad + polish pass 1 (best) & \textbf{0.1337} & \textbf{79.95\%} \\
\quad + passes 2, 3, 4 & 0.1369, 0.1439, 0.1449 & 79.41\%, 79.73\%, 79.15\% \\
PTQ, uniform + polish pass 1 & 0.1374 & 79.06\% \\
\bottomrule
\end{tabular}
\end{table}

\textbf{Findings.} (1) One zero-data polish pass nearly halves the
divergence ($0.2760\to0.1337$, $-51\%$) and lifts top-1 by $9.8$
points ($70.18\%\to79.95\%$): training only the quantization parameters
recovers half of the PTQ gap without a single external token. (2)
Hessian rank allocation beats uniform rank at \emph{identical} storage,
both before ($0.2760$ vs.\ $0.2951$, $-6.5\%$) and after ($0.1337$
vs.\ $0.1374$) polishing---the allocator's effect survives training
rather than being washed out by it. (3) Warm-start passes 2--4 degrade
monotonically ($0.1337\to0.1449$): each pass restarts the cosine
schedule at full learning rate, re-noising an already-converged point,
and the 512-row training pool shows a $\approx 0.05$ nats/token
generalization gap to the gate. Both effects point at on-policy
recovery (student-generated pools, reverse KL) as the precedented next
step \citep{wang2026onlineqat,agarwal2024gkd}, together with keeping
the best held-out checkpoint instead of the last. (4) For context, not
comparison: $0.134$ remains $2.7\times$ above our \emph{self-chosen}
$0.05$ gate---no literature-established 2-bit KL threshold exists; SLQ
reports distribution-level fidelity only at $5$--$6$ bits with
calibration data \citep{helcig2026slq}, and QAD's $0.004$ is a
4-bit-with-data number \citep{xin2026qad}. We make no near-lossless
claim. (5) These are attention-only gates on one 1.7B model;
likelihood-style metrics can also dissociate from autoregressive
generation quality for rotation-based methods \citep{prejzner2026bielik},
so a generation-coherence check joins the protocol before any stronger
claim.

\subsection{Full-Model Scaling on MiniCPM3-4B: Zero-Shot PTQ and the Exponential Perplexity Horizon}
\label{sec:minicpm3}

\textbf{Setting.} To evaluate whether 4D Clifford rotors scale to full-network parameterization beyond attention projections, we quantize all 248 heavy linear operators (\texttt{q\_a\_proj}, \texttt{q\_b\_proj}, \texttt{kv\_a\_proj}, \texttt{kv\_b\_proj}, \texttt{o\_proj}, \texttt{gate\_proj}, \texttt{up\_proj}, \texttt{down\_proj}) across all 62 layers of \texttt{openbmb/MiniCPM3-4B} (3.8B active parameters, 73,440-token vocabulary). MiniCPM3 features DeepSeek-style Multi-Head Latent Attention (MLA), where attention keys and queries are compressed into low-rank latent dimensions ($d_{k,\text{lora}}=256$, $d_{q,\text{lora}}=768$). Crucially, its head dimension $d_h=64=4^3$ aligns identically with three stages of 4D Clifford rotor butterflies ($\log_4(64)=3$), requiring zero padding and enabling complete SIMD lane utilization.

We evaluate two uniform quantization arms on the 65,536-token zero-data gate ($64 \times 1024$ teacher-generated sequences):
\begin{enumerate}[leftmargin=*]
\item \textbf{Arm A ($g=32$, uniform rank $r=8$):} Group scale size $g=32$, yielding an average information density of $\mathbf{2.119}$ bits per weight across all 248 linear modules.
\item \textbf{Arm B ($g=64$, uniform rank $r=8$):} Group scale size $g=64$, compressing the network down into the deep sub-2-bit regime at $\mathbf{1.859}$ bits per weight.
\end{enumerate}
Critically, both checkpoints represent \emph{pure zero-shot post-training quantization} (PTQ): weights were projected into the Clifford basis and ternary/SVD grids without any gradient backpropagation or QAT fine-tuning passes.

\begin{table}[t]
\centering\small
\caption{Full-model quantization results on \texttt{openbmb/MiniCPM3-4B} across all 248 linear projections (65,536-token zero-data teacher gate; forward $\mathrm{KL}(P_{\mathrm{fp16}}\|P_{\mathrm{quant}})$ in nats/token; vocabulary $|V|=73{,}440$). Both arms reflect pure zero-shot PTQ prior to gradient-based QAT polish.}
\label{tab:minicpm3}
\begin{tabular}{lccccc}
\toprule
Configuration & Group $g$ & bpw & KL (nats/token) & Top-1 Agreement & Student PPL (Teacher: 2.602) \\
\midrule
Teacher (BF16 baseline) & -- & 16.00 & 0.0000 & 100.00\% & 2.602 \\
Track 1 Arm A (Uniform) & 32 & 2.119 & \textbf{1.0473} & \textbf{76.45\%} & \textbf{6.798} \\
Track 1 Arm B (Uniform) & 64 & 1.859 & 1.2468 & 75.41\% & 8.465 \\
\bottomrule
\end{tabular}
\end{table}

\textbf{The Exponential Perplexity Breakdown.}
A striking feature of the empirical measurements in Table~\ref{tab:minicpm3} is the perplexity gap between the teacher ($2.602$) and student ($6.798$ on Arm A, $8.465$ on Arm B). We note that perplexity is fundamentally an exponential function of cross-entropy:
\begin{equation}
\label{eq:ppl_kl}
\mathrm{PPL}_{\mathrm{student}} = \exp\left( H(P_{\mathrm{teacher}}) + D_{\mathrm{KL}}(P_{\mathrm{teacher}} \parallel P_{\mathrm{student}}) \right) = \mathrm{PPL}_{\mathrm{teacher}} \cdot \exp\left( D_{\mathrm{KL}} \right).
\end{equation}
Equation~\eqref{eq:ppl_kl} demonstrates that forward KL divergence acts \emph{multiplicatively} on perplexity:
\begin{itemize}[leftmargin=*]
\item For Arm B ($g=64$, $D_{\mathrm{KL}} = 1.2468$ nats/token): $\exp(1.2468) = \mathbf{3.479\times}$ multiplier, predicting an expected student perplexity of $2.602 \times 3.479 = \mathbf{9.052}$ (closely matching the empirical measurement of $8.465$).
\item For Arm A ($g=32$, $D_{\mathrm{KL}} = 1.0473$ nats/token): $\exp(1.0473) = \mathbf{2.850\times}$ multiplier, predicting an expected student perplexity of $2.602 \times 2.850 = \mathbf{7.416}$ (empirically $6.798$).
\end{itemize}
Hence, the $+4.2$ to $+5.8$ PPL jump is not an architectural collapse, but the exact mathematical consequence of exponentiating the $\sim 1.05\text{ nat}$ discrete ternary entropy gap under zero-shot scalar quantization.

\textbf{Autoregressive Coherence and Decoding Attractors.}
Over an autoregressive generation of length $L$, sequence distribution fidelity scales as $\exp(-L \cdot D_{\mathrm{KL}}) \approx (0.75)^L$. Under strict greedy argmax decoding ($T=0$), small localized perturbations in output logits ($\Delta z \sim 10^{-2}$) can cause the model to fall into discrete attractor cycles (e.g., token repetition). However, under the Bonsai-2 sampling protocol ($T=0.7$, Top-$p=0.80$, presence penalty $1.5$), sampling directly traces the underlying 75.4\% teacher probability distribution, bypassing discrete argmax traps and sustaining syntactically and semantically fluent output on Apple Silicon Metal at $25.5$ tokens/second.

\textbf{Significance for Zero-Shot PTQ.}
In the published literature on sub-2-bit compression, standard zero-shot PTQ methods (such as naive RTN or uncalibrated GPTQ) experience catastrophic divergence ($\mathrm{PPL} > 1000$ or immediate newline loops) due to outlier-induced scale collapse. QuATernaryons achieves \textbf{76.45\% Top-1 token agreement} with the FP16 teacher over a 73,440-token vocabulary purely via algebraic geometric initialization. As demonstrated in recent representation fine-tuning literature \citep{malinovskii2024pvtuning,zhang2026catq}, refining only the continuous parameters $(\alpha, \beta, U, V)$ while keeping the discrete ternary grid $T$ and Clifford rotors $R$ frozen provides a direct gradient highway that collapses this zero-shot exponential penalty down to near-teacher parity.

\section{Related Work}
\label{sec:related}

\textbf{Weight PTQ below 4 bits.} GPTQ established Hessian-aware rounding
\citep{frantar2023gptq}; QuIP introduced incoherence processing with guarantees
\citep{chee2023quip}; QuIP\# combined the RHT with E8 lattice codebooks
\citep{tseng2024quipsharp}; AQLM used learned additive codebooks
\citep{egiazarian2024aqlm}; QTIP replaced codebooks with trellis-coded
quantization for ultra-high effective dimensionality \citep{tseng2024qtip};
NanoQuant reaches binary and sub-1-bit weights via low-rank binary
factorization \citep{chong2026nanoquant}; and differentiable bit-widths
co-optimize pruning with quantization down to $1.61$ bits
\citep{kang2026diffbitwidth}.
All treat the transform as an unstructured or Hadamard matrix; QuATernaryons
is orthogonal to codebook choice and can serve as the incoherence stage of any
of them. \textbf{Rotations for quantization.} QuaRot inserts fixed random
rotations end-to-end \citep{ashkboos2024quarot}; SpinQuant learns them on the
Stiefel manifold \citep{liu2024spinquant}; DuQuant adds permutation
\citep{lin2024duquant}; GSR reorders by sequency \citep{choi2025gsr}; HARP learns sparse
butterfly-orthogonal two-sided rotations initialized to the RHT and reports its
clearest gains at 2 bits \citep{zagitov2026harp}; ButterflyQuant learns
layer-adaptive two-wide Givens-butterfly transforms in place of the fixed
Hadamard rotation \citep{xu2025butterflyquant}; ParoQuant uses a scaled
pairwise rotation that is learnable yet cheap enough to apply online
\citep{liang2026paroquant}; and WUSH derives closed-form optimal linear
blockwise transforms for joint weight--activation quantization by combining a
Hadamard backbone with a data-dependent second-moment term
\citep{chen2026wush}. Our permutation--rotor
composition is structurally a learned, quaternion-parametrized 4-wide member of
this family (our two-wide learned-Givens baseline follows its
parameterization); the Great-Inversion survey \citep{jokar2026inversion} situates the
transform stage itself, and the QuATernaryons moment-equalizer fit is a concrete
calibrated point in the flattened pole it describes. \textbf{KV cache.} KIVI established per-channel/per-token 2-bit
asymmetry \citep{liu2024kivi}; KVQuant handles pre-attention outliers
\citep{hooper2024kvquant}; RotateKV applies outlier-aware adaptive rotations
for 2-bit KV via FWHT-compatible channel reordering \citep{su2025rotatekv};
KVTuner tunes mixed precisions \citep{li2025kvtuner}; QServe co-designs
W4A8KV4 serving \citep{lin2024qserve}; KV-COBRA co-optimizes per-head bit--rank
allocation under a fused Hadamard rotation \citep{ha2026kvcobra}; TaSQ tailors
the vector-quantization space via query-guided channel weighting and
covariance-aware grouping for 1-bit KV \citep{cheong2026tasq}.
\textbf{Geometric rotor quantization (2026).} TurboQuant frames online KV
vector quantization as dense-rotation + Lloyd--Max scalar quantization +
residual correction \citep{google2025turboquant}; RotorQuant (Pope) replaces
the dense rotation with blockwise $\mathrm{Cl}(3,0)$ rotors
\citep{pope2026rotorquant}; IsoQuant (Ji) upgrades the blocks to $4$D
isoclinic quaternion pairs, reports $4.5\times$ kernel speedups over
RotorQuant at matched MSE, and---importantly---validates only the stage-1
quantize--dequantize path on \emph{synthetic normalized i.i.d.\ vectors}, with
end-to-end model evaluation stated as future work \citep{ji2026isoquant}.
QuATernaryons extends this lineage in five directions: (a) the
\emph{weight side} at sub-2 bits (IsoQuant targets KV stage-1 only); (b)
\emph{multi-stage permutation--rotor composition} achieving global mixing---our
Table~\ref{tab:ablation} shows single-stage block-diagonal rotors, as in
IsoQuant, cannot disperse localized massive-activation spikes (error 1.094 vs.\
1.105 identity), a failure mode invisible in IsoQuant's i.i.d.-vector
evaluation; (c) \emph{closed-form calibration inside the exact rotor group}:
Lemma~\ref{lem:proc} gives an exact Rotor-Procrustes fit and
Theorem~\ref{thm:cayley} guarantees the double rotor realizes any block
target, whereas IsoQuant uses random Haar quaternions without a fitting theory;
and unlike WUSH's non-orthogonal closed-form transforms \citep{chen2026wush},
our optimum stays inside the rotor group, so exact isometry and inversion hold
by construction; (d) an explicit
attention-path perturbation bound (Prop.~\ref{prop:kv}); and (e) MSL/Triton
kernel specifications alongside their CUDA results. A distinct quaternion use
appears in HQMQ, which quantizes KV \emph{directions} to the 24-element Hurwitz
group composed with a random per-head codebook---a calibration-free quaternion
codebook rather than an orthogonal transform \citep{swain2026hqmq}; its
observation that random codebook seeds move end-task perplexity by less than
$1.5\%$, because left multiplication is an $S^3$ isometry, corroborates the
isometry-based argument used here. Independently, RotaryQuant's KV axis combines a Walsh--Hadamard stage with block-diagonal
$\mathrm{SO}(4)$ rotations for 3-bit KV inside a MoE serving system
\citep{lui2026rotaryquant}, confirming the viability of the block-rotor
operating point in production; QuATernaryons differs in calibrated quaternion
parametrization, sub-2-bit weight-side scope, and closed-form fitting theory.
We consider
\citep{ji2026isoquant} strictly prior and independent work on the shared
block-rotor idea, and position this paper as its systems-scale and
calibration-theoretic extension. \textbf{Clifford methods in ML.} Clifford-group
equivariant networks use rotors for geometric equivariance, not compression
\citep{ruhe2023clifford}. \textbf{Native low-bit training.}
BitNet b1.58 and OneBit address the QAT regime
\citep{wang2025bitnet,xu2024onebit}, complementary to our PTQ setting.
\textbf{QAT and self-distillation below 4 bits.} LLM-QAT established
data-free distillation from a model's own generations but only down to
4 bits \citep{liu2024llmqat}. TSLD already demonstrated distillation
QAT on \emph{ternary} weights in 2023, trained on real text (PTB) at
GPT-2/OPT scale \citep{kim2023tsld}. BitDistiller showed
same-architecture self-distillation can beat a larger teacher at 2-bit
\emph{asymmetric integer} weights \citep{du2024bitdistiller};
EfficientQAT's E2E-QP phase---freeze the quantized weights, train only
the quantization parameters---is the direct precedent of our polish
stage \citep{chen2025efficientqat}; mixed-domain calibration with
teacher-guided rectification and truncated forward-KL reaches 2-bit
Qwen3-8B \citep{okoshi2026reasoningqat}; and OnlineQAT moves recovery
onto \emph{student}-generated prefixes with reverse KL
\citep{wang2026onlineqat}, formalized by on-policy distillation
\citep{agarwal2024gkd}. NVIDIA's QAD reports KL $0.004$ to BF16 at
NVFP4 after distillation recovery \citep{xin2026qad}, while SLQ shows
distribution-level fidelity still requires $5$--$6$ bits even with
calibration data \citep{helcig2026slq}. To our knowledge, no prior work
combines \emph{zero external data} with a \emph{ternary} alphabet and a
distribution-fidelity gate; Section~\ref{sec:klgates} reports the first
such gates. \textbf{Low-rank repair and budgeted allocation.} Low-rank
residuals repair quantized weights \citep{saxena2025resq} or activations
\citep{scetbon2024lrc}; LR-QAT trains low-rank quantization-aware
parameters \citep{bondarenko2024lrqat}; BALANCER greedily allocates a
joint bit/rank budget across a model \citep{fan2026balancer}; QTEA
repairs a ternary base with salient $1{:}4$-sparse columns at $1.7$
effective bits \citep{guo2026qtea}. Our contribution here is the
ternary instantiation: Hessian-sensitivity \emph{rank} allocation with
a budget-matched uniform-rank ablation (Section~\ref{sec:klgates}).
\textbf{Kernels.} Marlin and successors show memory-bound dequant-fused GEMM
efficiency \citep{frantar2024marlin}; our Appendix~\ref{app:msl}--\ref{app:triton}
follow the same design grammar with 4-wide quaternion butterflies.

\section{Limitations}
\label{sec:limitations}

(1) \emph{Dilution gap.} With $T$ stages the uncalibrated diffusion factor is
$4^{T}$, not $d$; the RHT diffuses fully with $\log_2 d$ cheap butterflies. Our
value proposition is calibrated adaptivity at comparable cost, not asymptotic
dominance of the RHT (Prop.~\ref{prop:incoh}(ii)). (2) \emph{Scale of evidence.} Our experiments comprise controlled synthetic
ensembles, a toy attention model, a weight-side pilot on real trained
weights (Table~\ref{tab:real}), and distribution-fidelity gates on the
attention projections of one 1.7B model (Section~\ref{sec:klgates});
end-to-end perplexity, downstream evaluation, MLP/embedding coverage,
generation-coherence checks, and larger models remain future work
(Appendix~\ref{app:acc} specifies the protocol; we claim nothing beyond
it). The gates miss our own $0.05$ target by $2.7\times$; no
near-lossless claim is made.
(3) \emph{Kernel maturity.} The MSL/Triton listings in the appendices are
production-oriented specifications that have not yet been benchmarked on
silicon; their cost model is analytic. (4) \emph{Calibration cost.} The
coordinate-descent fit is offline and per-matrix, like all learned-rotation
methods; streaming KV calibration inherits RotateKV's head-wise statistics
assumption, which may interact with sink-token mitigation strategies
\citep{su2025rotatekv,sun2024massive}. (5) \emph{Grouping interaction.} Rotor
blocks pair channels at distance $d/16$ under the coarse stride; for models
whose outlier channels cluster at scales below 4, a finer first stage
($s_1=4$) is required and was observed to be slightly worse on our ensembles.
(6) \emph{Metric proxy.} Our weight-side evidence is reconstruction error,
which is a proxy rather than the objective: Zhao et al.\ show that a weight
matrix with lower calibration reconstruction loss can produce \emph{higher}
downstream loss once the activation distribution shifts, and refine the integer
codes within a fixed grid to be distributionally robust \citep{zhao2026drq}.
Because orthogonal transforms act on the input distribution in a controlled
way, we expect the effect to be milder here than for non-orthogonal transforms,
but we do not test it; applying such post-hoc refinement to rotor-quantized
codes adds no inference overhead and is a direct next step.

\section{Conclusion}

We introduced a quantization transform family built from 4D isoclinic Clifford
rotors: exactly orthogonal by algebra, 4-float parametrized per block,
inverted by conjugation, and composable into $O(d\log d)$ butterfly networks.
The theory pins down what rotors can and cannot do---coverage of $\SO(4)$
block-wise, parity of incoherence, a quantified dilution gap versus the
RHT---and the experiments show that with moment-equalizer calibration and
quantization-aware refinement the family dominates the random Hadamard
transform on every controlled ensemble while reducing dual-sided 2-bit KV
attention error by $3\times$ in a spiked toy model. On real trained 27B
weights, the calibrated transform reduces 2-bit reconstruction error by
$7.0\%$ over the identity and surpasses the random Hadamard transform on the
per-head KV projections. We release the
verification suite, full proofs, and kernel specifications, and we specify the
protocol for model-scale validation as the immediate next step.

\bibliographystyle{plainnat}
\bibliography{references}

\appendix

\section{Metal Shading Language kernel (rotor butterfly, float32 path)}
\label{app:msl}

The kernel applies a calibrated rotor stage to a row of $K$, $V$, or an
activation: one threadgroup transforms $4\times16$ channels per dispatch
tile; quaternions are loaded as \texttt{float4}; each 4-channel butterfly is
12 FLOP-multiply-adds with SIMD-4 ops. Dequantized 2-bit inputs use the same
butterfly (fused after dequant), mirroring Marlin-style fused dequant-GEMM
designs \citep{frantar2024marlin}.

\begin{lstlisting}[language=MSL,caption={Rotary-butterfly stage (left-rotor variant), MSL 3.1.}]
#include <metal_stdlib>
using namespace metal;

// One thread transforms one 4-channel block of one (row, head).
// qL: unit quaternion (float4) for this block; x: 4 input channels.
kernel void rotor_butterfly_f32(
    device const float4* x_in      [[buffer(0)]],   // [nRows*blocks]
    device const float4* qL        [[buffer(1)]],   // [blocks] unit quats
    device       float4* x_out     [[buffer(2)]],   // [nRows*blocks]
    constant uint& nRows           [[buffer(3)]],
    constant uint& blocks          [[buffer(4)]],
    uint2 tid [[thread_position_in_grid]])         // (row, blockIdx)
{
    const uint r = tid.x, b = tid.y;
    float4 q = qL[b];
    // Channel gather (stride permutation) is performed by the dispatch
    // wrapper into x_in; each thread sees its own permuted 4-channel group.
    float4 v = x_in[r*blocks + b];
    // Left-isoclinic butterfly: y = A(L) x (Definition 1); 12 FMA via
    // 4 SIMD-4 multiply-adds. Sign pattern verified against the paper's
    // matrix in verify_isoquant_math.py (left_matrix).
    float4 l = q;
    float4 vn = float4(-v.y, v.x,  v.w, -v.z);
    float4 vq = float4(-v.z,-v.w, v.x,  v.y);
    float4 vp = float4(-v.w, v.z,-v.y,  v.x);
    float4 y = l.xxxx*v + l.yyyy*vn + l.zzzz*vq + l.wwww*vp;
    x_out[r*blocks + b] = y;
}
// Inverse stage: same kernel with qL <- conj(qL) and inverse permutation;
// conj is 3 sign flips, so inversion is exact and free.
\end{lstlisting}

\begin{lstlisting}[language=MSL,caption={Fused 2-bit dequant + rotor butterfly for the KV path.}]
kernel void dequant2_rotor_f16(
    device const uchar*  packed  [[buffer(0)]], // 2-bit codes, [rows*ch/4]
    device const half2*  scales  [[buffer(1)]], // per group (scale, zero)
    device const float4* qL      [[buffer(2)]],
    device       float4* out     [[buffer(3)]],
    constant uint& nRows         [[buffer(4)]],
    constant uint& blocks        [[buffer(5)]],
    uint2 tid [[thread_position_in_grid]])
{
    const uint r = tid.x, b = tid.y;
    // unpack four 2-bit codes (two bytes) -> float4 after affine dequant
    const float s = float(scales[b].x), z = float(scales[b].y);
    const uchar c8  = packed[r*(blocks*2) + b];
    const uchar c8b = packed[r*(blocks*2) + b + blocks];
    float4 v = float4(s*float(c8  & 0x3)        - z,
                      s*float((c8  >> 2) & 0x3) - z,
                      s*float(c8b & 0x3)        - z,
                      s*float((c8b >> 2) & 0x3) - z);
    float4 l = qL[b];
    float4 vn = float4(-v.y, v.x,  v.w, -v.z);
    float4 vq = float4(-v.z,-v.w, v.x,  v.y);
    float4 vp = float4(-v.w, v.z,-v.y,  v.x);
    out[r*blocks + b] = l.xxxx*v + l.yyyy*vn + l.zzzz*vq + l.wwww*vp;
}
\end{lstlisting}

\section{Triton kernel specification (rotor butterfly)}
\label{app:triton}

\begin{lstlisting}[language=py,caption={Triton: per-head rotor stage fused with 2-bit dequant; BLOCK=4-aligned.}]
import triton, triton.language as tl

@triton.jit
def rotor_butterfly_kv2bit(
    X_ptr,      # *i8   packed 2-bit K or V, row-major [T, d]
    S_ptr, Z_ptr, # *fp32 per-group scale / zero [d/G]
    Q_ptr,      # *fp32 unit quaternions [d//4, 4]
    P_ptr,      # *i32  stride permutation [d]
    Y_ptr,      # *fp32 output [T, d]
    T, D, G: tl.constexpr, BLOCK_T: tl.constexpr):
    pid_t = tl.program_id(0)
    offs_t = pid_t * BLOCK_T + tl.arange(0, BLOCK_T)      # rows
    offs_d = tl.arange(0, D)                               # channels
    mask_t = offs_t < T
    # 1) dequant groupwise (levels {0,1,2,3} -> affine)
    codes = tl.load(X_ptr + offs_t[:, None] * D + offs_d[None, :],
                    mask=mask_t[:, None], other=0)
    g = offs_d // G
    s = tl.load(S_ptr + g); z = tl.load(Z_ptr + g)
    x = s[None, :] * codes.to(tl.float32) - z[None, :]
    # 2) 4-wide rotor butterflies, log4(D)+2 stages, stage tables in registers
    for t in tl.static_range(6):
        x = tl.reshape(x, (BLOCK_T, D // 4, 4))
        q0 = tl.load(Q_ptr + (tl.arange(0, D // 4) * 4 + 0))[:, None]
        q1 = tl.load(Q_ptr + (tl.arange(0, D // 4) * 4 + 1))[:, None]
        q2 = tl.load(Q_ptr + (tl.arange(0, D // 4) * 4 + 2))[:, None]
        q3 = tl.load(Q_ptr + (tl.arange(0, D // 4) * 4 + 3))[:, None]
        a, b, c, e = tl.split(x, 4)                        # components
        y0 = q0 * a - q1 * b - q2 * c - q3 * e
        y1 = q1 * a + q0 * b - q3 * c + q2 * e
        y2 = q2 * a + q3 * b + q0 * c - q1 * e
        y3 = q3 * a - q2 * b + q1 * c + q0 * e
        x = tl.reshape(tl.join(y0, y1, y2, y3), (BLOCK_T, D))
        # stride permutation between stages (double-buffer via index map)
        p = tl.load(P_ptr + offs_d)
        x = tl.gather(x, p[None, :].broadcast_to(BLOCK_T, D), axis=1)
    tl.store(Y_ptr + offs_t[:, None] * D + offs_d[None, :], x,
             mask=mask_t[:, None])
# Exact inverse: same kernel, Q <- conj(Q), P <- inverse permutation.
# Cost model: T*6 stages, each D/4 butterflies * 12 FMA + permutation pass
# ~= 18*D FLOP/row and 6*D bytes/row shuffled -> bandwidth-bound like FWHT.
\end{lstlisting}

\section{Accuracy protocol for model-scale validation, and cost accounting}
\label{app:acc}

\textbf{Protocol (pre-registered).} Models: Llama-2-7B/13B, Llama-3.1-8B.
Weights: 2-bit groupwise ($g=32$) with codebooks fixed to (a) uniform
symmetric and (b) E8-lattice; transforms compared: identity, RHT
\citep{tseng2024quipsharp}, learned Givens-block butterfly (2-wide), and the
calibrated rotor butterfly. Metrics: WikiText-2 perplexity (sequence 2048),
5-shot MMLU, GSM8K; KV: 2-bit with per-channel keys/per-token values
\citep{liu2024kivi}, long-context Needle-in-a-Haystack at 32k/128k, and
GSM8K-CoT following RotateKV's protocol \citep{su2025rotatekv} so results are
directly comparable to published numbers (RotateKV reports $<0.3$ PPL
degradation at 2-bit on Llama-2-13B).
Success criterion: rotor butterfly within 0.05 PPL of the best
transform at each bit-width, and strictly better on at least one long-context
metric.

\textbf{Cost accounting (analytic).} Per token, per head
($d_h=128$, $T=6$ stages): transform $=6\times32$ butterflies $\times12$ FMA
$=2{,}304$ FMA vs.\ FWHT's $7\times64\times4=1{,}792$ FMA ($1.29\times$);
parameter storage $6\times32\times8$ floats $=1{,}536$ floats/head vs.\ 0 for
fixed RHT and $128^2$ for dense Stiefel ($2.7\%$ thereof); exact inversion is
3 sign flips per quaternion. The 2-bit KV footprint with 128-group scales is
$0.25$ bytes/element $+$ scales, i.e.\ $8\times$ reduction vs.\ FP16,
identical to KIVI-class schemes; the transform adds no persistent per-token
state. Nominal bit-widths reference field widths rather than serialized bytes,
and ternary storage studies show the two diverge substantially in practice
\citep{georganas2026bitcos}.

\textbf{Reproducibility.} \texttt{verify\_isoquant\_math.py} (numpy-only)
regenerates every number in
Tables~\ref{tab:main}--\ref{tab:ablation} deterministically (seed 20260923)
and prints the 22 named checks; the paper's claims are stated so that each
check falsifies exactly one numbered assertion.

\end{document}
